% GENERATED FILE. Edit canonical lessons, then run npm run generate:books.
% canonical-source-hash: fnv1a64:6a3aee9aced6ff31
% canonical-lessons: PA-C142-vinga, PA-C142-camkila, PA-C142-caura, PA-C142-bhira, PA-C142-gilla

\chapter{Crooked, Shiny, Wide, Narrow, Wet}
\label{ch:pa-crooked-shiny-wide-narrow-wet}
\hlchaptermodality{142}

\begin{chapteropening}
\textbf{By the end of this chapter:} \emph{I can say crooked, shiny, wide, narrow and wet.}

Complete the last lesson of chapter 142: I can say crooked, shiny, wide, narrow and wet.
\end{chapteropening}

\section[\texorpdfstring{\pa{ਵਿੰਗਾ} (vingā)}{ਵਿੰਗਾ (vingā)}]{\pa{ਵਿੰਗਾ} (vingā) — crooked}

\label{lesson:PA-C142-vinga}



\begin{quote}
Before the new one: say the Punjabi for round, then the Punjabi for straight.
\end{quote}

\subsection*{You'll want to know: \pa{ਵਿੰਗਾ}}
\textbf{\pa{ਵਿੰਗਾ}} — \emph{vingā} — \textquotedblleft{}crooked\textquotedblright{}.

\begin{grammarlens}[title={What you've built}]
One more word for this chapter.
\end{grammarlens}

\subsection*{Your turn}
\begin{itemize}
\raggedright
  \item \emph{Say it:} \emph{vingā}
  \item \emph{Say it:} \emph{vingā}, once more
  \item \emph{Say it:} say \emph{vingā}
  \item \emph{Recall:} say the Punjabi for round, then the Punjabi for straight, then say \emph{vingā} again
\end{itemize}

\begin{tcolorbox}[breakable,colback=teal!4,colframe=teal!35!black,title={Before you move on}]
What does \pa{ਵਿੰਗਾ} mean? (\textquotedblleft{}Crooked\textquotedblright{}.) Say it once more.
\end{tcolorbox}
\section[\texorpdfstring{\pa{ਚਮਕੀਲਾ} (camkīlā)}{ਚਮਕੀਲਾ (camkīlā)}]{\pa{ਚਮਕੀਲਾ} (camkīlā) — shiny}

\label{lesson:PA-C142-camkila}



\begin{quote}
Before the new one: say the Punjabi for straight, then the Punjabi for crooked.
\end{quote}

\subsection*{You'll want to know: \pa{ਚਮਕੀਲਾ}}
\textbf{\pa{ਚਮਕੀਲਾ}} — \emph{camkīlā} — \textquotedblleft{}shiny\textquotedblright{}.

\begin{grammarlens}[title={What you've built}]
One more word for this chapter.
\end{grammarlens}

\subsection*{Your turn}
\begin{itemize}
\raggedright
  \item \emph{Say it:} \emph{camkīlā}
  \item \emph{Say it:} \emph{camkīlā}, once more
  \item \emph{Say it:} say \emph{camkīlā}
  \item \emph{Recall:} say the Punjabi for straight, then the Punjabi for crooked, then say \emph{camkīlā} again
\end{itemize}

\begin{tcolorbox}[breakable,colback=teal!4,colframe=teal!35!black,title={Before you move on}]
What does \pa{ਚਮਕੀਲਾ} mean? (\textquotedblleft{}Shiny\textquotedblright{}.) Say it once more.
\end{tcolorbox}
\section[\texorpdfstring{\pa{ਚੌੜਾ} (cauṛā)}{ਚੌੜਾ (cauṛā)}]{\pa{ਚੌੜਾ} (cauṛā) — wide}

\label{lesson:PA-C142-caura}



\begin{quote}
Before the new one: say the Punjabi for crooked, then the Punjabi for shiny.
\end{quote}

\subsection*{You'll want to know: \pa{ਚੌੜਾ}}
\textbf{\pa{ਚੌੜਾ}} — \emph{cauṛā} — \textquotedblleft{}wide\textquotedblright{}.

\begin{grammarlens}[title={What you've built}]
One more word for this chapter.
\end{grammarlens}

\subsection*{Your turn}
\begin{itemize}
\raggedright
  \item \emph{Say it:} \emph{cauṛā}
  \item \emph{Say it:} \emph{cauṛā}, once more
  \item \emph{Say it:} say \emph{cauṛā}
  \item \emph{Recall:} say the Punjabi for crooked, then the Punjabi for shiny, then say \emph{cauṛā} again
\end{itemize}

\begin{tcolorbox}[breakable,colback=teal!4,colframe=teal!35!black,title={Before you move on}]
What does \pa{ਚੌੜਾ} mean? (\textquotedblleft{}Wide\textquotedblright{}.) Say it once more.
\end{tcolorbox}
\section[\texorpdfstring{\pa{ਭੀੜਾ} (bhīṛā)}{ਭੀੜਾ (bhīṛā)}]{\pa{ਭੀੜਾ} (bhīṛā) — narrow}

\label{lesson:PA-C142-bhira}



\begin{quote}
Before the new one: say the Punjabi for shiny, then the Punjabi for wide.
\end{quote}

\subsection*{You'll want to know: \pa{ਭੀੜਾ}}
\textbf{\pa{ਭੀੜਾ}} — \emph{bhīṛā} — \textquotedblleft{}narrow\textquotedblright{}.

\begin{grammarlens}[title={What you've built}]
One more word for this chapter.
\end{grammarlens}

\subsection*{Your turn}
\begin{itemize}
\raggedright
  \item \emph{Say it:} \emph{bhīṛā}
  \item \emph{Say it:} \emph{bhīṛā}, once more
  \item \emph{Say it:} say \emph{bhīṛā}
  \item \emph{Recall:} say the Punjabi for shiny, then the Punjabi for wide, then say \emph{bhīṛā} again
\end{itemize}

\begin{tcolorbox}[breakable,colback=teal!4,colframe=teal!35!black,title={Before you move on}]
What does \pa{ਭੀੜਾ} mean? (\textquotedblleft{}Narrow\textquotedblright{}.) Say it once more.
\end{tcolorbox}
\section[\texorpdfstring{\pa{ਗਿੱਲਾ} (gillā)}{ਗਿੱਲਾ (gillā)}]{\pa{ਗਿੱਲਾ} (gillā) — wet}

\label{lesson:PA-C142-gilla}



\begin{quote}
Before the new one: say the Punjabi for wide, then the Punjabi for narrow.
\end{quote}

\subsection*{You'll want to know: \pa{ਗਿੱਲਾ}}
\textbf{\pa{ਗਿੱਲਾ}} — \emph{gillā} — \textquotedblleft{}wet\textquotedblright{}.

\begin{grammarlens}[title={What you've built}]
One more word for this chapter.
\end{grammarlens}

\subsection*{Your turn}
\begin{itemize}
\raggedright
  \item \emph{Say it:} \emph{gillā}
  \item \emph{Say it:} \emph{gillā}, once more
  \item \emph{Say it:} say \emph{gillā}
  \item \emph{Recall:} say the Punjabi for wide, then the Punjabi for narrow, then say \emph{gillā} again
\end{itemize}

\begin{tcolorbox}[breakable,colback=teal!4,colframe=teal!35!black,title={Before you move on}]
What does \pa{ਗਿੱਲਾ} mean? (\textquotedblleft{}Wet\textquotedblright{}.) Say it once more.
\end{tcolorbox}
